\section{Архитектура системы}
\label{sec:architecture}


\subsection{Общая схема пайплайна}


В качестве базовой платформы выбран YouTube по двум причинам: это крупнейший агрегатор видеоконтента и образовательного контента в частности,
охватывающий материалы от университетских курсов до авторских лекций на любом языке и по любой дисциплине;
и для него существует развитая экосистема инструментов извлечения данных — в частности, yt-dlp \cite{ytdlp2021},
обеспечивающий доступ к аудио, видео, субтитрам и метаданным в рамках публичного API платформы.
Архитектурно пайплайн не привязан к YouTube: замена источника данных требует адаптации только первого слоя сбора,
тогда как все последующие модули работают с видеофайлом, аудиодорожкой и транскриптом — форматами, не специфичными для конкретной платформы.
Это включает и локальные файлы при необходимости офлайн-обработки.

Система построена как четырёхслойный каскад, в котором каждый слой принимает результаты предыдущего и передаёт обогащённое представление следующему.

\textbf{Слой 1 — сбор данных.} На вход подаётся ссылка на YouTube-видео. Метаданные, аудиопоток и видеокадры извлекаются через yt-dlp и ffmpeg \cite{ffmpeg2026};
транскрипт — из субтитров YouTube либо через Whisper \cite{radford2022whisper} при их отсутствии.

\textbf{Слой 2 — извлечение признаков.} Параллельно по трём каналам: из аудио извлекаются просодические и акустические характеристики;
из видеокадров — метрики качества изображения, характеристики сцены и поведения спикера; из транскрипта — текстовые и семантические признаки.

\textbf{Слой 3 — анализ.} На основе извлечённых признаков вычисляются метрики и флаги: техническое качество, флаги восприятия, профиль подачи.
На этом же слое выполняется LLM-анализ содержания. Результаты слоя 3 интерпретируются с учётом контекста пользователя.

\textbf{Слой 4 — синтез и выдача.} Формируется итоговая аналитическая выдача:
таймкоды с нарративной суммаризацией, численные метрики, флаги и агрегированная метрика соответствия запросу.

\subsection{Контекст пользователя}


Система принимает четыре параметра пользовательского профиля:

\begin{itemize}
\item \textbf{Кто я} (USER\_TYPE) — роль пользователя: обучающийся, автор контента, образовательная организация или заказчик контента.
\item \textbf{Уровень образования} (LEVEL) — от школьника до специалиста / профессионала; для авторов контента — уровень целевой аудитории.
\item \textbf{Уровень погружённости} (IMMERSION\_LEVEL) — от «изучаю с нуля» до «экспертный уровень».
\item \textbf{Цель просмотра} (VIEW\_GOAL) — составить общее представление, закрыть точечные вопросы или последовательно изучить тему.
\end{itemize}

Выбор именно этих четырёх параметров продиктован стремлением к минимальному достаточному набору:
вместе они покрывают ключевые сценарии использования и не требуют от пользователя избыточных усилий.
Контекст пользователя пронизывает систему на нескольких уровнях: гранулярность тематической сегментации определяется целью просмотра;
нарратив формируется не как универсальное описание видео, а как объяснение его ценности для конкретного профиля;
интерпретация профиля подачи зависит от уровня и погружённости пользователя; метрика соответствия запросу агрегируется с учётом полного профиля.
Для авторов контента состав выдачи расширяется: помимо содержательных таймкодов формируются таймкоды проблемных зон с привязкой к конкретным флагам.

\subsection{Состав аналитической выдачи}


Система формирует пять типов выходных данных:

\begin{itemize}
\item \textbf{Техническое качество} — агрегированная оценка аудио и видео по воспринимаемым характеристикам.
\item \textbf{Флаги восприятия} — бинарные сигналы, фиксирующие наличие паттернов подачи, негативно влияющих на восприятие.
\item \textbf{Профиль подачи} — положение видео на двух независимых осях: академичность и инструментальность.
\item \textbf{Нарративная суммаризация} — тематические сегменты с таймкодами и текстовое описание ценности видео относительно запроса пользователя.
\item \textbf{Метрика соответствия запросу} — интегральная оценка, агрегирующая все предыдущие компоненты с учётом пользовательского профиля.
\end{itemize}
